\documentclass[11pt]{article}
\usepackage[a4paper,margin=1in]{geometry}
\usepackage[T1]{fontenc}
\usepackage[utf8]{inputenc}
\usepackage{lmodern}
\usepackage{amsmath,amssymb,amsthm}
\usepackage{microtype}
\usepackage{hyperref}
\title{Emmy Noether\\English Translation Draft\\On Integral Rational Representation of Invariants}
\author{}
\date{}
\begin{document}
\maketitle

\section*{On integral rational representation of the invariants of a system of arbitrarily many basic forms}

At the end of his contribution to the Schwarz Festschrift, \emph{On the invariants of a system of arbitrarily many basic forms}, Hilbert expressed the conjecture that these invariants can be represented integrally and rationally by the finitely many invariants of the system \((J,PJ)\), where \(J\) denotes a complete system of invariants for the linearly independent basic forms and \(PJ\) denotes the invariants obtained from \(J\) by polar processes.

In what follows I show, in Part I, that this conjecture is indeed correct, and in fact as a simple consequence of the reduction theorem: every form in arbitrarily many rows of \(n\) variables can be represented as a sum of polars of forms that involve only \(n\) fixed rows.

In Part II I give a more conceptual proof of the reduction theorem by interpreting the question as a problem of equivalence of linear families of forms. I owe this point of view, and likewise an essential part of the reasoning employed, to Fischer's paper \emph{On algebraic systems of modules} and to unpublished propositions of E. Fischer, whose use he kindly permitted me. Finally, in Part III I indicate how analogous considerations also lead to the most general series expansions.

\section*{I. The invariant-theoretic consequence of the reduction theorem}

Let \(\Theta\) be a form homogeneous in each of the \(N>n\) rows
\[
A_1,A_2,\dots,A_N,
\]
where in general the row \(A_i\) consists of the \(n\) elements
\[
a_1^{(i)},\dots,a_n^{(i)}.
\]
The coefficients of \(\Theta\) may be indeterminates or quantities belonging to a given field of rationality. Let \(P\) denote an integral rational function---with rational numerical coefficients---of the polar operators
\[
P_{ih}
=
a_1^{(i)}\frac{\partial}{\partial a_1^{(h)}}
+
a_2^{(i)}\frac{\partial}{\partial a_2^{(h)}}
+\cdots+
a_n^{(i)}\frac{\partial}{\partial a_n^{(h)}},
\]
where by a product such as \(P_{ih}P_{jk}\Theta\) one understands the action of \(P_{ih}\) on \(P_{jk}\Theta\).

Then the reduction theorem referred to above is expressed by the identity
\[
\Theta = \sum P\Theta_0,
\tag{1}
\]
where the forms \(\Theta_0\) involve only the rows
\[
A_1,A_2,\dots,A_n
\]
and are obtained from \(\Theta\) by polar processes \(P\).

Now consider the basic system \((F)\) consisting of the \(N\) forms of the same order and the same number of variables,
\[
F_1,F_2,\dots,F_N,
\]
whose coefficients are taken respectively to be the \(N\) rows
\[
A_1,A_2,\dots,A_N.
\]
Let \((J)\) denote a complete system of invariants of
\[
F_1,\dots,F_n,
\]
that is, a finite set of invariants such that every invariant of
\[
F_1,\dots,F_n
\]
can be expressed integrally and rationally through the invariants \((J)\). Hilbert's conjecture to be proved is then the following: for the simultaneous invariants \(S\) of
\[
F_1,\dots,F_N
\]
a complete system is given by the invariants \((J,PJ)\).

Applying the reduction theorem \((1)\) to a simultaneous invariant \(S\)---which can itself be regarded as a form in the \(N\) rows \(A_i\)---one sees at once that every simultaneous invariant is representable as a sum of polars of invariants that involve only the \(n\) rows
\[
A_1,\dots,A_n.
\]
But such forms are invariants of
\[
F_1,\dots,F_n,
\]
and hence can be expressed integrally and rationally through the invariants \((J)\). On the other hand, by the very definition of the polar processes, every form obtained from an invariant of
\[
F_1,\dots,F_n
\]
by a polar process is again a simultaneous invariant of
\[
F_1,\dots,F_N.
\]
Since the polar process preserves integral rational expressibility, it follows that every simultaneous invariant \(S\) of
\[
F_1,\dots,F_N
\]
can be represented integrally and rationally through the invariants \((J,PJ)\).

Thus we have proved:

\medskip
\noindent\textbf{Theorem.}
Every simultaneous invariant of a system of arbitrarily many basic forms is integrally and rationally expressible through the finite family \((J,PJ)\), where \(J\) is a complete invariant system for \(n\) linearly independent basic forms and \(PJ\) denotes all invariants obtained from \(J\) by polar operations.

\medskip
This establishes Hilbert's conjecture.

\bigskip

\end{document}
